\cleardoublepage
\chapter*{\listofabbrevname}
\phantomsection
\addcontentsline{toc}{chapter}{\listofabbrevname}

\begin{acronym}[DSP48E1]

%%% Zkratky %%%
\acro{ARM}
    {Architektura procesoru (Advanced RISC Machine)}

\acro{BRAM}
    {Blok interní paměti FPGA (Block Random Access Memory)}

\acro{DC}
    {Stejnosměrný (Direct Current)}

\acro{DSP48E1}
    {Hardwarový blok násobiče v FPGA (Digital Signal Processing block)}

\acro{FF}
    {Klopný obvod (Flip-Flop)}

\acro{HDL}
    {Jazyk pro popis hardware (Hardware Description Language)}

\acro{HIL}
    {Hardware ve smyčce (Hardware in the Loop)}

\acro{IPMSM}
    {Synchronní motor s permanentními magnety uvnitř rotoru (Interior Permanent Magnet Synchronous Motor)}

\acro{IP}
    {Blok duševního vlastnictví pro FPGA (Intellectual Property core)}

\acro{LUT}
    {Vyhledávací tabulka (Lookup Table)}

\acro{MTPA}
    {Maximální moment na ampér (Maximum Torque Per Ampere)}

\acro{MTPV}
    {Maximální moment na volt (Maximum Torque Per Volt)}

\acro{PI}
    {Proporcionálně-integrační regulátor}

\acro{PL}
    {Programovatelná logika čipu Zynq (Programmable Logic)}

\acro{PMSM}
    {Synchronní motor s permanentními magnety (Permanent Magnet Synchronous Motor)}

\acro{PS}
    {Procesorový subsystém čipu Zynq (Processor Subsystem)}

\acro{SISO}
    {Systém s jedním vstupem a jedním výstupem (Single Input Single Output)}

\acro{SoC}
    {Systém na čipu (System on Chip)}

\acro{SVPWM}
    {Prostorově vektorová pulzně šířková modulace (Space Vector Pulse Width Modulation)}

\bigskip
\noindent\textbf{Symboly}
\bigskip

%%% Latinské symboly %%%

\acro{An}
    [\ensuremath{A_n}]
    {Matice výstupních napěťových vektorů SVPWM pro sektor $n$}

\acro{b}
    [\ensuremath{b}]
    {Koeficient viskózního tření \hfill [Nm\,s/rad]}

\acro{E}
    [\ensuremath{E}]
    {Matice křížové vazby v modelu motoru v souřadnicích $dq_{1,3}$}

\acro{G}
    [\ensuremath{G(p)}]
    {Přenosová funkce soustavy}

\acro{C}
    [\ensuremath{C(p)}]
    {Přenosová funkce regulátoru}

\acro{i}
    [\ensuremath{i}]
    {Statorový proud \hfill [A]}

\acro{idq}
    [\ensuremath{i_{dq_{13}}}]
    {Vektor statorového proudu v souřadnicích $dq_{1,3}$ \hfill [A]}

\acro{idqref}
    [\ensuremath{i_{dq_{13}}^*}]
    {Referenční (požadovaný) vektor statorového proudu \hfill [A]}

\acro{ihat1}
    [\ensuremath{\hat{I}_1}]
    {Amplituda 1.\,harmonické fázového proudu \hfill [A]}

\acro{ihat3}
    [\ensuremath{\hat{I}_3}]
    {Amplituda 3.\,harmonické fázového proudu \hfill [A]}

\acro{IDCmax}
    [\ensuremath{I_{DC\text{max}}}]
    {Maximální proud odebíraný z DC zdroje \hfill [A]}

\acro{iswmax}
    [\ensuremath{i_{sw\text{-max}}}]
    {Maximální proud tranzistorem \hfill [A]}

\acro{J}
    [\ensuremath{J}]
    {Moment setrvačnosti rotoru \hfill [kg\,m$^2$]}

\acro{Jopt}
    [\ensuremath{J_{min}}]
    {Nákladová funkce optimalizace MTPA \hfill [A$^2$]}

\acro{k}
    [\ensuremath{k}]
    {Konstanta příspěvku 3.\,harmonické permanentního magnetu \hfill [--]}

\acro{k3}
    [\ensuremath{k_3}]
    {Poměr amplitud 1.\,a 3.\,harmonické proudu (MTPV) \hfill [--]}

\acro{Kp}
    [\ensuremath{K_p}]
    {Proporcionální zesílení regulátoru \hfill [--]}

\acro{Ks0}
    [\ensuremath{K_{s0}}]
    {Clarkova transformační matice (včetně nulové složky) z $abcde$ do $\alpha\beta_{1,3}0$}

\acro{Kr0}
    [\ensuremath{K_{r0}(\theta)}]
    {Transformační matice z $abcde$ do rotujících souřadnic $dq_{1,3}0$}

\acro{Ks}
    [\ensuremath{K_s}]
    {Clarkova transformační matice (bez nulové složky) z $abcde$ do $\alpha\beta_{1,3}$}

\acro{Kr}
    [\ensuremath{K_r}]
    {Transformační matice (bez nulové složky) z $abcde$ do $dq_{1,3}$}

\acro{Las}
    [\ensuremath{L_{A5}}]
    {Konstanta symetrické složky vzájemné indukčnosti statoru \hfill [H]}

\acro{Lb}
    [\ensuremath{L_{B}}]
    {Konstanta pro poloze závislé části indukčnosti statoru \hfill [H]}

\acro{Ld1}
    [\ensuremath{L_{d1},\,L_{q1}}]
    {Indukčnosti statoru v souřadnicích $dq_1$ \hfill [H]}

\acro{Ld3}
    [\ensuremath{L_{d3},\,L_{q3}}]
    {Indukčnosti statoru v souřadnicích $dq_3$ \hfill [H]}

\acro{Lls}
    [\ensuremath{L_{ls}}]
    {Rozptylová indukčnost statorového vinutí \hfill [H]}

\acro{Lm}
    [\ensuremath{L_m}]
    {Fluktuace indukčnosti statoru \hfill [H]}

\acro{Ls}
    [\ensuremath{L_s(\theta)}]
    {Matice vzájemných indukčností statoru \hfill [H]}

\acro{Ma5}
    [\ensuremath{M_{a5}}]
    {Matice symetrické části vzájemných indukčností statorových vinutí}

\acro{Mb5}
    [\ensuremath{M_{b5}(\theta)}]
    {Matice poloze závislé části vzájemných indukčností statorových vinutí}

\acro{Ms}
    [\ensuremath{M_s}]
    {Vzájemná indukčnost statorových vinutí \hfill [H]}

\acro{n}
    [\ensuremath{n}]
    {Číslo sektoru SVPWM \hfill [--]}

\acro{p}
    [\ensuremath{p}]
    {Počet pólových dvojic \hfill [--]}

\acro{R0}
    [\ensuremath{R_0(\theta)}]
    {Rotační matice z $\alpha\beta_{1,3}0$ do $dq_{1,3}0$}

\acro{Rs}
    [\ensuremath{R_s}]
    {Odpor vinutí statoru na fázi \hfill [$\Omega$]}

\acro{swi}
    [\ensuremath{sw_i}]
    {Spínací signál tranzistoru $i$-té fáze střídače \hfill [--]}

\acro{Te}
    [\ensuremath{T_e}]
    {Elektrický (elektromagnetický) moment motoru \hfill [Nm]}

\acro{Teref}
    [\ensuremath{T^*}]
    {Požadovaný moment \hfill [Nm]}

\acro{Ti}
    [\ensuremath{T_i}]
    {Integrační časová konstanta regulátoru \hfill [s]}

\acro{TL}
    [\ensuremath{T_L}]
    {Zátěžový moment \hfill [Nm]}

\acro{TM}
    [\ensuremath{T_M}]
    {Vzorkovací perioda modulátoru \hfill [s]}

\acro{Tm}
    [\ensuremath{T_m}]
    {Mechanický moment motoru \hfill [Nm]}

\acro{Ts}
    [\ensuremath{T_s}]
    {Vzorkovací perioda řízení \hfill [s]}

\acro{v}
    [\ensuremath{v}]
    {Statorové napětí \hfill [V]}

\acro{Vdc}
    [\ensuremath{V_{dc}}]
    {Napájecí napětí DC meziobvodu \hfill [V]}

%%% Řecké symboly %%%

\acro{gamma}
    [\ensuremath{\gamma}]
    {Vektor příspěvku rotorového magnetického toku ke statorovým tokům}

\acro{theta}
    [\ensuremath{\theta}]
    {Elektrický úhel natočení rotoru \hfill [rad]}

\acro{thetar}
    [\ensuremath{\theta_r}]
    {Mechanický úhel natočení rotoru \hfill [rad]}

\acro{thetam}
    [\ensuremath{\theta_m}]
    {Měřený mechanický úhel polohy rotoru \hfill [rad]}

\acro{tau}
    [\ensuremath{\tau}]
    {Časová konstanta filtru žádané hodnoty momentu \hfill [s]}

\acro{omegafr}
    [\ensuremath{\omega_{fr}}]
    {Elektrická úhlová rychlost rotoru \hfill [rad/s]}

\acro{omegam}
    [\ensuremath{\omega_m}]
    {Mechanická úhlová rychlost \hfill [rad/s]}

\acro{omegamref}
    [\ensuremath{\omega_m^*}]
    {Požadované (referenční) otáčky \hfill [rad/s]}

\acro{omegae}
    [\ensuremath{\omega_e}]
    {Elektrická úhlová rychlost napájecího napětí \hfill [rad/s]}

\acro{omegabase}
    [\ensuremath{\omega_{base}}]
    {Základní (jmenovitá) rychlost \hfill [rad/s]}

\acro{omegac}
    [\ensuremath{\omega_c}]
    {Frekvence řezu regulátoru \hfill [rad/s]}

\acro{psi}
    [\ensuremath{\psi}]
    {Magnetický tok \hfill [Wb]}

\acro{PsiPM}
    [\ensuremath{\Psi_{PM}}]
    {Toková vazba permanentního magnetu \hfill [Wb]}

\acro{a5}
    [\ensuremath{a_5}]
    {Základní fázor pětifázové soustavy, $a_5 = e^{j2\pi/5}$ \hfill [--]}

\end{acronym}
